\documentclass[11pt]{article}
\usepackage[a4paper,margin=0.85in]{geometry}
\usepackage{amsmath,amssymb,amsthm,mathtools}
\usepackage{booktabs,longtable,array,tabularx}
\usepackage{enumitem}
\usepackage{microtype}
\usepackage{xcolor}
\usepackage{hyperref}
\usepackage{url}
\usepackage{listings}
\usepackage{graphicx}
\usepackage{float}
\usepackage{seqsplit}
\usepackage{fancyhdr}
\usepackage{titlesec}
\usepackage[T1]{fontenc}
\usepackage[utf8]{inputenc}

\hypersetup{colorlinks=true,linkcolor=blue,citecolor=blue,urlcolor=blue}
\pagestyle{fancy}
\fancyhf{}
\lhead{Forensic Formal Research Program}
\rhead{September 2026}
\cfoot{\thepage}
\setlength{\headheight}{14pt}

\newtheorem{theorem}{Theorem}[section]
\newtheorem{lemma}[theorem]{Lemma}
\newtheorem{proposition}[theorem]{Proposition}
\newtheorem{corollary}[theorem]{Corollary}
\newtheorem{definition}[theorem]{Definition}
\newtheorem{conjecture}[theorem]{Conjecture}
\newtheorem{obligation}[theorem]{Open Obligation}
\newtheorem{counterexample}[theorem]{Counterexample}
\newtheorem{remark}[theorem]{Remark}

\newcommand{\Status}[1]{\texttt{\detokenize{#1}}}
\newcommand{\R}{\mathbb{R}}
\newcommand{\C}{\mathbb{C}}
\newcommand{\N}{\mathbb{N}}
\newcommand{\eps}{\varepsilon}
\newcommand{\norm}[1]{\left\lVert #1\right\rVert}
\newcommand{\abs}[1]{\left\lvert #1\right\rvert}
\newcommand{\Repo}{\texttt{ShantiDraconis/navier-stokes-critical-barrier-audit}}

\title{\textbf{A Provenance-Aware Formal Research Program for Critical Barriers, Error Operators, and Computational Conjecture Auditing}\\
\large Mathematical laws, theorem obligations, reproducible tests, and a peer-review submission architecture}
\author{Tiago Paschoalatto Fagliari\\Research identity: ShantiDraconis}
\date{12 September 2026}

\begin{document}
\maketitle

\begin{abstract}
This article consolidates a multi-year computational and formal research corpus into an auditable mathematical program. The contribution is deliberately separated into four layers: (i) exact provenance and chronology, (ii) proved elementary identities and operator laws, (iii) computationally testable model behavior, and (iv) unresolved bridge theorems required before any classical Millennium-problem conclusion can be asserted. The provenance layer includes Google Colab metadata from 2024--2026, exact SHA-256 fingerprints of recovered notebooks, Git commit and blob histories, and a public audit repository. The formal layer introduces a generic defect--correction--recomputation operator, proves normalization, contraction, stability, scaling, conjugation, and non-inference lemmas, and states explicit open obligations for Navier--Stokes and Riemann-hypothesis-inspired branches. For three-dimensional incompressible Navier--Stokes, the diagnostic ratio
\[
\eps_{NS}(u)=\frac{\norm{(u\cdot\nabla)u}_{L^2}}{\nu\norm{\Delta u}_{L^2}}
\]
is scale invariant whenever defined, but this fact is not a regularity theorem. For the SE'ET/Riemann-inspired branch, the off-critical-line envelope
\[
S_d(u)=e^{(1/2+d)u}\sin(\gamma u)
\]
obeys exact amplitude and residual identities, but no zero of the Riemann zeta function is established by those identities. The paper therefore converts an exploratory corpus into a peer-reviewable research program whose claims are typed by proof status and whose remaining mathematical gaps are stated as explicit theorem obligations rather than silently assumed conclusions.
\end{abstract}

\section{Purpose, claim boundary, and contribution}
The project contains numerical experiments, symbolic models, proof-assistant source files, repository histories, and later forensic reconstructions. Such a corpus can be scientifically useful only if chronology, implementation, theoremhood, and conjecture are kept distinct. We therefore adopt the following rule.

\begin{definition}[Four-axis claim typing]
Every material claim is assigned four logically independent attributes:
\begin{enumerate}[label=(\roman*)]
\item \emph{provenance status}: what exact object existed, where, and when;
\item \emph{formal status}: definition, theorem statement, checked proof, assumption-dependent proof, sketch, computational evidence, or placeholder;
\item \emph{scientific scope}: what mathematical or physical statement the object actually entails;
\item \emph{priority scope}: repository-local first occurrence, corpus-local first occurrence, public publication date, or global historical priority.
\end{enumerate}
\end{definition}

The present work claims only what can be supported on those axes. In particular, neither Git chronology nor notebook age proves a Millennium Prize Problem, and no theorem below upgrades an unresolved bridge to a solved theorem by terminology alone. The official Clay problem statements remain the governing classical targets \cite{ClayNS,ClayRH}.

\section{Forensic chronology and exact digital objects}
\subsection{Recovered notebook objects}
Five notebook objects uploaded into the audit session were fingerprinted exactly. Their byte hashes are part of the reproducibility record:

\begin{longtable}{@{}p{0.28\textwidth}r p{0.55\textwidth}@{}}
\toprule Object & Bytes & SHA-256 \\ \midrule \endhead
\texttt{Untitled9 (1).ipynb} & 324 & \texttt{\seqsplit{298b16a180218da825d562bcbb29f3b26955074fea202ace7faadae051dd8632}}\\
\texttt{Untitled9.ipynb} & 369483 & \texttt{\seqsplit{ab41de3e75ba6eaf80a58474d689328777b915cb3e4321a79bdce9634892439d}}\\
\texttt{Untitled10.ipynb} & 339761 & \texttt{\seqsplit{373b8ca210135c5660e17ba534c176aaabbc71bb79082fe24b2492791c1eac8a}}\\
\texttt{Untitled12 (1).ipynb} & 233487 & \texttt{\seqsplit{72169dafb7b3f3811f8fd38a2fea77649975974caf0b10698d0b472fb59ff6c3}}\\
\texttt{Untitled0.ipynb} & 586384 & \texttt{\seqsplit{e308a48344053f781385b357c6ca5d626db5e07a2c83dea107970af3f599fdd3}}\\
\bottomrule
\end{longtable}

The exact-byte statement is strong: the listed hash identifies the uploaded byte sequence. The historical-date statement is weaker unless a historical Drive revision can be independently downloaded and hashed.

\begin{lemma}[Hash identity]
Let $H$ be a collision-resistant cryptographic hash treated operationally as an exact byte fingerprint for the audit. If two archived objects have identical byte sequences, then they have identical SHA-256 values.
\end{lemma}
\begin{proof}SHA-256 is deterministic. Equal inputs therefore yield equal outputs.\end{proof}

\begin{remark}The converse is treated as operationally sufficient for archival identity in this audit, but collision resistance is a security assumption, not a theorem of pure mathematics.\end{remark}

\subsection{Chronology currently supported}
The documented sequence includes user-supplied Google Drive version/detail records and directly recovered Git history:
\[
\begin{array}{rcl}
14\text{ Aug }2024 &:& \text{earliest displayed Untitled3 lineage version},\\
19\text{ Aug }2024 &:& \text{later Untitled3 versions},\\
15\text{ Sep }2024 &:& \text{a 7.1 MB Colab created and modified},\\
08\text{ Oct }2024 &:& \text{displayed Untitled9 current-version timestamp},\\
24\text{--}25\text{ Oct }2024 &:& \text{another Colab created, modified, and opened},\\
26\text{ Feb }2025 &:& \text{the Sep-2024 Colab reopened},\\
23\text{ Nov }2025 &:& \text{earliest directly recovered Navier--Stokes Git origin},\\
16\text{ Apr }2026 &:& \text{displayed Untitled0 Drive lineage versions},\\
12\text{ Sep }2026 &:& \text{exact current notebook-byte audit and consolidation}.
\end{array}
\]

\begin{theorem}[Chronology does not imply global priority]
Suppose object $A$ has an authenticated timestamp earlier than object $B$ in the audited corpus. Then one may infer that $A$ predates $B$ in that corpus. One may not infer from this fact alone that the mathematical idea encoded by $A$ was first discovered globally at that date.
\end{theorem}
\begin{proof}The timestamp relation compares only the observed corpus. Global priority quantifies over all earlier public and private work by all researchers. The corpus timestamp contains no information sufficient to exclude unobserved prior art.\end{proof}

\section{Generic defect--correction formalism}
\begin{definition}[State, reference, and defect]
Let $(X,d)$ be a metric or seminormed state space, let $x_*\in X$ be a reference state, and let $E:X\to[0,\infty]$ be a defect functional. A simple instance is $E(x)=d(x,x_*)$.
\end{definition}

\begin{definition}[Defect-dependent correction operator]
A correction architecture is a family $\mathcal C_\theta:X\times[0,\infty]\to X$, with update $\mathcal T_\theta(x)=\mathcal C_\theta(x,E(x))$ and recomputed defect $E^+(x)=E(\mathcal T_\theta x)$.
\end{definition}

\begin{definition}[Closed correction loop]
A sequence $(x_n)$ is a closed correction loop if
\[
D_n:=E(x_n),\qquad x_{n+1}=\mathcal C_{\theta_n}(x_n,D_n),\qquad D_{n+1}=E(x_{n+1}).
\]
A pipeline lacking the final recomputation $D_{n+1}$ is not, by this definition, a closed error-feedback loop.
\end{definition}

\begin{theorem}[Uniform contraction closes the abstract loop]
Assume there exists $q\in[0,1)$ such that $E(\mathcal T_\theta x)\le qE(x)$ for every admissible state $x$. Then $E(x_n)\le q^nE(x_0)$ and $E(x_n)\to0$.
\end{theorem}
\begin{proof}Induction gives the geometric bound; $q^n\to0$.\end{proof}

\begin{obligation}[Problem-specific coercivity]
For any Millennium-problem application, prove that the chosen defect $E$ controls a recognized classical quantity strongly enough that $E\to0$ or $\sup E<\infty$ implies the desired theorem.
\end{obligation}

\section{Weighted error operators and coordinate geometry}
Let $\Delta_i\ge0$, $g_i:[0,\infty)\to[0,\infty)$, $w_i\ge0$, $\sum_iw_i=1$, and
\[
E_w(\Delta)=\sum_{i=1}^n w_i g_i(\Delta_i).
\]
If $E_w>0$, define $C_i^{(w)}=100w_i g_i(\Delta_i)/E_w$.

\begin{lemma}[Percentage normalization]
If $E_w>0$, then $\sum_iC_i^{(w)}=100$.
\end{lemma}
\begin{proof}Immediate from the definition.\end{proof}

\begin{theorem}[Reparameterization non-canonicity]
The vector $(C_i^{(w)})$ is generally not invariant under nonlinear reparameterization.
\end{theorem}
\begin{proof}With $n=2$, equal weights, identity transforms, and $\Delta=(1,2)$, the percentages are $(100/3,200/3)$. Replacing the second coordinate by its square gives contributions proportional to $(1,4)$, hence $(20,80)$.\end{proof}

For $\eps(P)\in\R_+^n$, $\norm{\eps}_2=(\sum_i\eps_i^2)^{1/2}$.
\begin{lemma}$\norm{\eps}_2=0$ iff every $\eps_i=0$.\end{lemma}
\begin{proof}A finite sum of nonnegative squares is zero exactly when every square is zero.\end{proof}

\begin{counterexample}[Monotone error reduction need not solve the problem]
$E_n=1+1/(n+1)$ is strictly decreasing but converges to $1$, not $0$.
\end{counterexample}

\section{Navier--Stokes branch}
On $\R^3$ consider
\begin{align}
\partial_t u +(u\cdot\nabla)u+\nabla p &= \nu\Delta u+f,\\
\nabla\cdot u &=0.
\end{align}
Classical theory includes Leray, Caffarelli--Kohn--Nirenberg, and endpoint work of Escauriaza--Seregin--\v{S}ver\'ak \cite{Leray1934,CKN1982,ESS2003}; see also Temam \cite{Temam2001}.

Define, when the denominator is nonzero,
\[
\eps_{NS}(u)=\frac{\norm{(u\cdot\nabla)u}_{L^2}}{\nu\norm{\Delta u}_{L^2}}.
\]

\begin{proposition}[Scale invariance]
Under $u_\lambda(x,t)=\lambda u(\lambda x,\lambda^2t)$, $\eps_{NS}(u_\lambda)=\eps_{NS}(u)$ whenever both sides are defined.
\end{proposition}
\begin{proof}Both $(u_\lambda\cdot\nabla)u_\lambda$ and $\Delta u_\lambda$ scale pointwise by $\lambda^3$; the $L^2(\R^3)$ change of variables yields a norm factor $\lambda^{3/2}$ for both.\end{proof}

\begin{remark}Scaling invariance is structurally interesting but is not a regularity criterion.\end{remark}

\begin{obligation}[NS critical bridge]
Prove, under fully stated hypotheses,
\[
\mathcal H_{\mathrm{repo}}(u;0,T)\Longrightarrow u\in L^\infty(0,T;L^3(\R^3)),
\]
where $\mathcal H_{\mathrm{repo}}$ is the exact repository condition.
\end{obligation}

\begin{obligation}[Non-circularity]Show that $\mathcal H_{\mathrm{repo}}$ can be established from hypotheses weaker than the desired regularity conclusion.\end{obligation}
\begin{obligation}[Denominator and degeneracy]Specify rigorously what happens when $\norm{\Delta u}_2=0$.\end{obligation}
\begin{obligation}[Threshold coercivity]Any universal threshold requires an independent coercive inequality connecting it to a classical critical norm.\end{obligation}

\begin{lemma}[Integer-period mean of a pure mode]
For $\omega\ne0$ and integer $m\ge1$, the mean of $Ae^{i\omega t}$ over $2\pi m/\omega$ is zero.
\end{lemma}
\begin{counterexample}[Zero mean is not pointwise control]$\sin t$ has zero period mean but $L^\infty$ norm $1$.\end{counterexample}

\section{Riemann/SE'ET branch}
The Riemann hypothesis concerns nontrivial zeros of $\zeta(s)$ and asserts real part $1/2$ \cite{Titchmarsh1986,Edwards1974,IwaniecKowalski2004,ClayRH}.

Define
\[
R(u)=e^{u/2}\sin(\gamma u),\qquad S_d(u)=e^{(1/2+d)u}\sin(\gamma u).
\]
\begin{lemma}[Exact residual factorization]
\[
S_d(u)-R(u)=e^{u/2}(e^{du}-1)\sin(\gamma u).
\]
\end{lemma}
\begin{proof}Factor the common term.\end{proof}

\begin{lemma}[Envelope ratio]
Whenever $\sin(\gamma u)\ne0$, $S_d(u)/R(u)=e^{du}$.
\end{lemma}
For $d=0.025$, $u=15$, $e^{du}=e^{0.375}\approx1.45499$.

Define $\rho_{\sigma,\tau}=1/2+\sigma d+i\tau\gamma$, $\sigma,\tau\in\{-1,+1\}$.
\begin{lemma}[Conjugation closure]The quartet is closed under complex conjugation.\end{lemma}
\begin{lemma}[Critical-line reflection]The map $s\mapsto1-\overline{s}$ maps $\rho_{\sigma,\tau}$ to $\rho_{-\sigma,\tau}$.\end{lemma}

These identities do not prove zerohood.

\begin{obligation}[Zeta-zero authentication]
For every purported zeta zero $\rho$, independently establish $\zeta(\rho)=0$ with a certified procedure and rigorous error bound.
\end{obligation}
\begin{obligation}[Explicit-formula hypotheses]
State the exact explicit formula, smoothing/truncation convention, multiplicities, remainder and domain.
\end{obligation}

Define
\[
K_d(\gamma)=\left(\frac{\gamma}{2\pi}\right)^{2d},\quad
w_L=\frac{K_d}{1+K_d},\quad w_R=\frac1{1+K_d}.
\]
\begin{lemma}[Weight normalization]$w_L+w_R=1$.\end{lemma}
\begin{proposition}[Monotonicity in $\gamma$]For $d>0$, $K_d(\gamma)$ is strictly increasing for $\gamma>0$.\end{proposition}
\begin{proof}$K_d'(\gamma)=2d(2\pi)^{-2d}\gamma^{2d-1}>0$.\end{proof}

For $\Delta K(u_{\max})=e^{du_{\max}}-e^{-du_{\max}}=2\sinh(du_{\max})$:
\begin{lemma}[Gamma independence by definition]$\Delta K$ is independent of $\gamma$.\end{lemma}

\section{Earlier computational branches and rigorous reinterpretation}
The 2024 corpus contains exploratory constants, singularities, particle/molecular visualizations, SU-labeled models, quantum circuits, signal processing, weather models and chronology. Physical labels do not establish physical theories.

\begin{definition}[Dimensional homogeneity]A physical equation is dimensionally homogeneous when all additive terms share the same physical dimension and both sides have the same dimension.\end{definition}
\begin{theorem}[Inhomogeneous sums are not physical laws without normalization]
If $a_1,\ldots,a_n$ have incompatible physical dimensions, the bare sum $a_1+\cdots+a_n$ is not a dimensionally meaningful physical observable unless an explicit nondimensionalization or unit-converting map is supplied.
\end{theorem}

\begin{lemma}[Reciprocal divergence]For $p>0$, $x^{-p}\to+\infty$ as $x\to0^+$.\end{lemma}
\begin{remark}This does not define a new algebraic value for $0/0$.\end{remark}

\section{Formal proof-assistant discipline}
Formal verification should follow proof-assistant standards \cite{Lean2015,IsabelleHOL2002,Norell2009}.
\begin{definition}[Kernel-checked theorem status]
A theorem receives \Status{FORMAL_PROOF_CHECKED} only if the exact source revision is identified, prover/dependencies are pinned, clean compilation succeeds, no admitted placeholder lies in the transitive dependency chain, and all nontrivial axioms are enumerated.
\end{definition}
A field such as \texttt{deepTheorem : Prop} is a specification obligation until an instance supplies a proof term.

\section{Test suite specification}
Required tests include exact notebook hashes and cell counts; percentage normalization; zero-vector characterization; contraction; the monotone-nonzero counterexample; SE'ET residual and envelope identities; stiffness monotonicity; weight normalization; gamma-independence of the defined $\Delta K$; quartet symmetries; NS scaling; zero-mean counterexample; and anti-overclaim checks for timestamps, zeta-zero labels, placeholders and hard-coded residuals.

\section{Peer-review graph and evidence nodes}
Recommended nodes are \Status{DRIVE_OBJECT}, \Status{RAW_NOTEBOOK}, \Status{CELL_OCCURRENCE}, \Status{GIT_COMMIT}, \Status{GIT_BLOB}, \Status{FORMAL_THEOREM}, \Status{COMPUTATION}, \Status{CLAIM}, \Status{PUBLICATION}, and \Status{BRIDGE}. Edges should be typed as \Status{EXACT_COPY}, \Status{DERIVED_FROM}, \Status{RENAMED}, \Status{REFINED}, \Status{GENERALIZED}, \Status{SPECIALIZED}, \Status{SYNTACTIC_TRANSLATION}, \Status{CONCEPTUAL}, \Status{CONTRADICTS}, or \Status{UNRESOLVED}.

\section{Authorship and provenance claims that survive review}
Strong bounded language is: ``the audited corpus contains an account-associated object dated $T$, and an exact recovered revision/blob contains formula $F$ with hash $H$.'' Corpus-priority language is: ``among all objects inspected under the declared search procedure, this is the earliest verified occurrence of $F$.'' Global-first claims require a separate prior-art study.

\section{Gap ledger}
\begin{longtable}{@{}p{0.04\textwidth}p{0.31\textwidth}p{0.49\textwidth}p{0.1\textwidth}@{}}
\toprule ID & Gap & Closure condition & Status \\ \midrule \endhead
G1 & Drive revision identity & Download historical revision; SHA-256; bind cell to timestamp & OPEN\\
G2 & 2024--2025 continuity & Find exact intermediate notebook/repository objects & OPEN\\
G3 & First occurrence of $0/0$, SU, $C_I$, SBFE, collapse & Per-concept revision/commit ledger & OPEN\\
G4 & Exact notebook-to-Git derivation & Exact or normalized cell/blob matches plus chronology & OPEN\\
G5 & Multi-prover clean build & Pin tools/dependencies; compile without admitted dependencies & OPEN\\
G6 & NS repository condition $\Rightarrow L^\infty_tL^3_x$ & Complete non-circular proof & OPEN\\
G7 & NS threshold coercivity & Prove quantitative estimate & OPEN\\
G8 & SE'ET parameter $\rho$ is zeta zero & Certified zeta evaluation/zero isolation & OPEN\\
G9 & Explicit-formula completeness & State exact formula with truncation/error terms & OPEN\\
G10 & High-precision integrity & Eliminate float downcasts; document precision end-to-end & OPEN\\
G11 & Random experiment reproducibility & Seed RNGs; archive data; confidence intervals & OPEN\\
G12 & Physical dimensional consistency & Nondimensionalize or remove physical-law interpretation & OPEN\\
G13 & Global priority & Exhaustive external prior-art study & OPEN\\
G14 & Independent review & External theorem/provenance review and responses & OPEN\\
\bottomrule
\end{longtable}

\section{Submission strategy}
The corpus should be decomposed by claim type. A paper centered on experiments, exact model identities, counterexamples, and sharply stated conjectures is aligned with \emph{Experimental Mathematics} \cite{ExperimentalMathScope}. After clean multi-prover builds and a genuine formal-method contribution, a venue such as the \emph{Journal of Automated Reasoning} becomes appropriate. A specialist Navier--Stokes paper should be submitted only after G6/G7 are closed. The reproducibility software can be separated as a research-software article, for example for \emph{SoftwareX} \cite{SoftwareXScope}.

Recommended sequence: freeze cryptographic provenance; publish software/audit infrastructure; submit the experimental-mathematics manuscript; complete multi-prover formalization; submit specialist NS/RH theorem papers only if the classical bridges are actually discharged.

\section{Conclusion}
The strongest established results are provenance facts, exact computational identities, elementary operator theorems, scale behavior, and falsification tests for invalid inference patterns. The most important unresolved results are a non-circular Navier--Stokes critical-space bridge and certified zeta-zero linkage for the Riemann-inspired branch. Exposing those gaps as theorem obligations makes the project more defensible.

\section*{Data and code availability}
The public audit repository is \path{ShantiDraconis/navier-stokes-critical-barrier-audit} \cite{RepoAudit}. Related repositories include \path{ShantiDraconis/universal}, \path{ShantiDraconis/millennium-navier-stokes-I}, and \path{ShantiDraconis/Trans} \cite{RepoUniversal,RepoNS,RepoTrans}.

\section*{Claim-scope declaration}
No statement in this manuscript should be read as a claim that the Navier--Stokes or Riemann Millennium Prize Problems have been solved.

\begin{thebibliography}{99}
\bibitem{Leray1934} J. Leray, ``Sur le mouvement d'un liquide visqueux emplissant l'espace,'' \emph{Acta Mathematica}, 63 (1934), 193--248.
\bibitem{CKN1982} L. Caffarelli, R. Kohn, and L. Nirenberg, ``Partial regularity of suitable weak solutions of the Navier--Stokes equations,'' \emph{Communications on Pure and Applied Mathematics}, 35(6) (1982), 771--831.
\bibitem{ESS2003} L. Escauriaza, G. Seregin, and V. \v{S}ver\'ak, ``$L_{3,\infty}$-solutions of Navier--Stokes equations and backward uniqueness,'' \emph{Russian Mathematical Surveys}, 58(2) (2003), 211--250.
\bibitem{Temam2001} R. Temam, \emph{Navier--Stokes Equations: Theory and Numerical Analysis}, AMS Chelsea, 2001.
\bibitem{Titchmarsh1986} E. C. Titchmarsh, revised by D. R. Heath-Brown, \emph{The Theory of the Riemann Zeta-Function}, 2nd ed., Oxford University Press, 1986.
\bibitem{Edwards1974} H. M. Edwards, \emph{Riemann's Zeta Function}, Academic Press, 1974.
\bibitem{IwaniecKowalski2004} H. Iwaniec and E. Kowalski, \emph{Analytic Number Theory}, AMS Colloquium Publications 53, 2004.
\bibitem{Lean2015} L. de Moura et al., ``The Lean Theorem Prover (System Description),'' \emph{CADE-25}, LNCS 9195, Springer, 2015, 378--388.
\bibitem{IsabelleHOL2002} T. Nipkow, L. C. Paulson, and M. Wenzel, \emph{Isabelle/HOL: A Proof Assistant for Higher-Order Logic}, LNCS 2283, Springer, 2002.
\bibitem{Norell2009} U. Norell, ``Dependently Typed Programming in Agda,'' \emph{Advanced Functional Programming}, LNCS 5832, Springer, 2009, 230--266.
\bibitem{Sandve2013} G. K. Sandve et al., ``Ten Simple Rules for Reproducible Computational Research,'' \emph{PLoS Computational Biology}, 9(10) (2013), e1003285.
\bibitem{ClayNS} Clay Mathematics Institute, ``Navier--Stokes Equation: Millennium Prize Problem,'' official problem description, accessed 12 September 2026.
\bibitem{ClayRH} Clay Mathematics Institute, ``Riemann Hypothesis: Millennium Prize Problem,'' official problem description, accessed 12 September 2026.
\bibitem{ExperimentalMathScope} Taylor \& Francis, ``Experimental Mathematics: Aims and Scope,'' journal website, accessed 12 September 2026.
\bibitem{SoftwareXScope} Elsevier, ``SoftwareX: Research Software Journal,'' journal information, accessed 12 September 2026.
\bibitem{RepoAudit} T. Paschoalatto Fagliari, \emph{navier-stokes-critical-barrier-audit}, GitHub repository, 2026.
\bibitem{RepoUniversal} T. Paschoalatto Fagliari, \emph{universal}, GitHub repository, 2025--2026.
\bibitem{RepoNS} T. Paschoalatto Fagliari, \emph{millennium-navier-stokes-I}, GitHub repository, 2025--2026.
\bibitem{RepoTrans} T. Paschoalatto Fagliari, \emph{Trans}, GitHub repository, archival recovery branch, 2026.
\end{thebibliography}
\end{document}
